
\section{Interface Design}\fbeg{}\begin{frame}{\ft{Interface Design Principles for the Research Object Context}}


\vspace{-.5em}	

%\MyDiamond{}
{\Large\fontfamily{uhv}\selectfont
%\begin{center}
\hspace*{2pt}\begin{minipage}{.96\textwidth}
\vspace{4pt}

%\definecolor{blback}{RGB}{0,100,100}	
%\definecolor{blfront}{RGB}{0,100,50}

%\fcolorbox{lqboutercolor}{lqbinnercolor}{\begin{minipage}{\textwidth}%
		
\begin{lightquadblockc}{0,0.1,0.1}{\parbox{21cm}{\centering \vspace{10pt}What are some design patterns that should be encouraged?\vspace{3pt}}}
\begin{center}\begin{minipage}{1.03\textwidth}
{\LARGE \setlength{\leftmargini}{30pt}\begin{enumerate}
\dmitem \textbf{Unify Pass-by-Reference Initialization and Pointer Return} \hspace{.5em} Called procedures 
agnostic as to the source of initialized memory buffers/objects (whether 
pre-initialized reference arguments, delegating ownership of heap-allocated memory, or Return Value Optimization) 

{\vsp}

\dmitem \textbf{Programmer Control of Variable Lifetimes}  \hspace{.5em} Lifetimes 
can be associated with different lexical scopes of other variables'

{\vsp}

\dmitem \textbf{\q{Overload by Contract}}  \hspace{.5em}  Programmers supply 
different versions of a procedure depending on whether contract-tests 
are known at compile time to have been evaluated 

{\vsp}

\dmitem \textbf{Easy Foreign Function Interface}  \hspace{.5em} Calling 
native-compiled procedures should be possible with simple registration fuctions

{\vsp}

\dmitem \textbf{Self-Documenting}  \hspace{.5em} Given $y = f(x)$, 
is $y$ being initialized or re-initialized?  Does $f$ have side-effects? 
Was $y$ default constructed?  
%Is $x$ passed by reference or by value?Const

{\vsp}

\dmitem We 
can guarantee $const$.  Can we \textsl{guarantee} to modify 
as with an lvalue? 
 



\end{enumerate}
}\end{minipage}
\end{center}
\end{lightquadblockc}
\end{minipage}

%}
%\end{minipage}
%\end{center}
}

\end{frame}

